\chapter{MARCO METODOLÓGICO}

La metodología adapta el esquema de compuertas incrementales de los trabajos previos a un
sistema que no se entrena de punta a punta. Cada capacidad de CAVA se valida con una
compuerta pre-registrada: el criterio, los datos y el umbral se fijan antes de ejecutar, y el
split de test de cada dominio se consulta una sola vez por hipótesis. El pipeline es
reproducible: cada corrida guarda un manifiesto, los experimentos se registran en orden
cronológico (EXP-009 a EXP-040), las decisiones de arquitectura se documentan como ADR y una
suite de pruebas automatizadas cubre los módulos centrales. El Anexo~C resume los
experimentos citados.

\section{Datos y splits}

\begin{itemize}
  \item MovieNet \cite{huang2020movienet}, dominio película. Las anotaciones de escena con
        la convención de BaSSL \cite{mun2022bassl} se usan para construir una referencia
        aproximada del mapa de tópicos: dos tomas pertenecen al mismo tópico si comparten
        reparto (proxy de reparto). Los splits de desarrollo (10 películas), confirmación y
        test (47 películas) son disjuntos y el de test no se usa para ajustar nada.
  \item BBC Planet Earth: 11 episodios documentales de $\sim$49 minutos con anotaciones de
        escena de 5 anotadores. Valida el perfil de documental contra el acuerdo humano.
  \item AMI \cite{carletta2007ami}: 120 reuniones de $\sim$35 minutos con anotación de
        tópicos y de hablantes. Valida el perfil de reunión y la identidad de hablante.
  \item VidChapters \cite{ventura2024vidchapters}, dominio video web: blind40 (40 videos de
        10 a 40 minutos) para la evaluación ciega contra capítulos de autor, y raw\_ext (50
        videos) para desarrollo.
  \item TVSum \cite{song2015tvsum} y SumMe \cite{gygli2014summe}: video de usuario de 1 a 11
        minutos con importancia anotada por personas, usada como referencia automática de
        la calidad de los destacados.
  \item VoxConverse \cite{chung2020voxconverse} y AVA-Speech: conjuntos de audio para la
        identidad de hablante y los filtros de habla.
  \item Conjunto A/B de EXP-040: 9 videos educativos de YouTube (divulgación, comentario,
        conversación y tutorial; 5 en español y 4 en inglés; de 21 a 57 minutos) sin
        anotación previa, reunidos para la evaluación con personas.
\end{itemize}

Ninguna corrida del producto usa anotaciones en la inferencia: el contrato de corrida
\texttt{cava-run-v1} registra \texttt{gt\_used: false} en el manifiesto de cada una (115
corridas).

\section{Pipeline experimental}

\begin{enumerate}
  \item Unidades y señales: ffprobe y ffmpeg preparan el audio y las unidades fijas;
        faster-whisper large-v3 transcribe; SigLIP codifica un frame por unidad; se extraen
        log-mel y tono; el VLM (Qwen2.5-VL-7B) describe bloques de $\sim$40\,s.
  \item Gestor: la función \texttt{manage()} se ejecuta con los hiperparámetros del perfil
        (percentil de corte, quórum, tolerancia, modo de enlace, percentil de retorno y
        abstención) y las fronteras se ajustan a los cortes visuales certificados por
        SegLoop.
  \item Mapa y etiquetas: \texttt{topic\_map.py} produce los dos niveles, los retornos y la
        procedencia; Qwen2.5-14B etiqueta los tópicos; los nombres de personajes se propagan
        por los enlaces.
  \item Versiones: \texttt{versions.py} genera los cinco formatos con clips mp4, listas de
        reproducción e índice, y G-versión verifica los presupuestos por video.
  \item Evaluación: un script por experimento lee los JSON de \texttt{results/}, calcula
        los criterios pre-registrados y el veredicto se anota en el registro de
        experimentos. Las recompensas de clips se calculan en un archivo aparte
        (\texttt{rewards.jsonl}) para que la evaluación no modifique la corrida.
\end{enumerate}

\section{Protocolo estadístico}

\begin{itemize}
  \item Pre-registro: cada compuerta declara criterio, dirección, conjuntos y umbral antes
        de abrir los datos de evaluación, con fecha en el registro de experimentos. El
        umbral típico es $z > 2$, donde $z$ es la diferencia media pareada por video
        dividida por su error estándar.
  \item Splits: el de desarrollo se usa para diagnosticar y calibrar, el de confirmación
        para decidir y el de test solo para la hipótesis final. EXP-018 confirmó en test (47
        películas) lo diagnosticado en desarrollo y decidido en confirmación.
  \item Comparaciones fuera del pre-registro: se reportan como análisis exploratorio y no
        como validación. Una confirmación previa sobre el split de entrenamiento quedó
        marcada como post hoc.
  \item Revisiones con juez: los revisores no saben de qué brazo viene cada caso, los
        controles se emparejan por película y desfase, y la significación se calcula con la
        prueba exacta de Fisher. El primer lote de G-enlace tuvo revisores distintos por
        brazo; se identificó el problema y se lanzó un segundo lote con el protocolo
        corregido.
\end{itemize}

\section{Atribución de la ganancia del gestor}

El gestor de CAVA difiere de la base \texttt{fused} en dos aspectos a la vez: no combina las
señales y, en su configuración final, usa enlace promedio con abstención. La comparación
pre-registrada en EXP-018 (\texttt{cava\_abst} frente a \texttt{fused}) mide el efecto
conjunto. Para separar los dos aspectos se reportan también, sobre las mismas 47 películas,
las comparaciones de las variantes sin abstención (\texttt{cava}, con componentes conexas, y
\texttt{cava\_avg}, con enlace promedio) frente a \texttt{fused}, y de \texttt{cava\_abst}
frente a \texttt{cava\_avg}. Estas comparaciones se calcularon con el mismo script que la
principal y se presentan como análisis secundario. La abstención no tiene equivalente en la
base fusionada, porque al concatenar los embeddings deja de existir una señal individual que
pueda abstenerse.

\section{Sistema de compuertas}

\begin{table}[H]
    \centering
    \caption{Compuertas de validación de CAVA y su estado al congelar el sistema.}
    \label{tab:gates}
    \begin{tabular}{p{2.6cm}p{6.6cm}p{3.6cm}}
        \hline
        \textbf{Compuerta} & \textbf{Criterio} & \textbf{Estado} \\
        \hline
        Test del gestor & MCC(\texttt{cava\_abst}) $-$ MCC(\texttt{fused}) $> 0$ con
                      $z > 2$ en las 47 películas de test & Cumplida ($z = 2.46$) \\
        G-coherencia & Margen intra$-$inter positivo en cada señal y en cada película
                      & Cumplida (47/47 películas en las tres señales) \\
        G-versión & Toda lista existe, cumple su presupuesto y hay al menos una versión
                      por perfil & Cumplida (58/58, 80/80, 90/90, 25/25, 9/9) \\
        G-enlace-2 & Enlaces no contiguos juzgados como el mismo hilo más que momentos
                      lejanos al azar (revisión ciega, Fisher $p < 0.05$) & No cumplida
                      ($p = 0.72$) \\
        Calidad de clips & Destacados y recapitulación superan a bases simples en
                      referencias automáticas & No cumplida (EXP-031/037/038) \\
        Narrador (EXP-039) & Un detector con habla filtrada detecta al hablante único
                      mejor que el actual ($z > 2$) sin más falsos & No cumplida
                      ($z = -0.37$ y $0.00$) \\
        EXP-040 (A/B) & Preferencia humana de las versiones frente a bases propias &
                      Diseñada; pre-registro pendiente \\
        \hline
    \end{tabular}
\end{table}

\section{Infraestructura y reproducibilidad}

Los experimentos se ejecutan en el clúster institucional (GPU NVIDIA RTX A6000 de 48\,GB y
A100 vía SLURM) con caché local de HuggingFace. El VLM y los modelos de voz corren en GPU y el
gestor en CPU. Cada corrida del producto escribe un manifiesto (commit de git, indicador de
cambios sin confirmar, perfil e hiperparámetros) y el archivo \texttt{clips.jsonl}. La suite
de 25 pruebas automatizadas cubre el gestor, el agregador adaptativo, el ajuste a cortes, el
mapa y la compuerta de versiones, y se ejecuta antes de cada envío a la cola. El paquete
\texttt{topicagg} declara sus dependencias en \texttt{pyproject.toml}. La decisión de no
fusionar representaciones y el uso de votantes activos con abstención están documentados en
la ADR-0001.

\section{Métricas}

\begin{itemize}
  \item Mapa de tópicos: MCC sobre la clasificación de pares de unidades (mismo tópico o
        no) contra la referencia (proxy de reparto en MovieNet; anotación de tópicos en
        AMI). El número de tópicos y la tasa de enlace (fracción de pares enlazados) se usan
        para diagnosticar sobresegmentación y percolación. La precisión de enlace se compara
        con la tasa base de pares positivos.
  \item Coherencia: margen entre la similitud coseno dentro de cada grupo (dejando una unidad
        fuera) y la similitud entre grupos, por señal, con la fracción de películas en que es
        positivo.
  \item Identidad de hablante: ARI contra la anotación (AMI, VoxConverse) y especificidad de
        las identidades publicadas.
  \item Fronteras: F1 contra cortes de referencia (VidChapters, BBC) y distancia del ajuste a
        los cortes visuales certificados.
  \item Calidad de clips: recompensa igual al $z$ de la fracción de anotadores que eligió cada
        segundo (TVSum, SumMe), comparada con bases de clips equiespaciados, clip central y
        azar.
  \item Versiones: existencia y cumplimiento del presupuesto de cada lista (G-versión), por
        video y formato.
\end{itemize}
